\chapter[Eliptyczne zagadnienia brzegowe]{Eliptyczne zagadnienia brzegowe: formy dywergencyjne, nierówność Gårdinga i alternatywa Fredholma}
\label{chap:pde13}

W rozdziałach~\ref{chap:pde-variational} i~\ref{chap:pde12} zbudowaliśmy
słabą teorię laplasjanu. Teraz współczynniki będą zależały od punktu,
możliwe będą człony pierwszego i zerowego rzędu, a forma nie musi być
symetryczna. Rozdzielimy trzy pytania: czy operator jest eliptyczny,
czy związana z nim forma jest koercywna, i czy odpowiednie zagadnienie
brzegowe ma dokładnie jedno rozwiązanie. Wykazanie samej eliptyczności
nie rozstrzyga ostatniego pytania.

\section{Postać dywergencyjna a niedywergencyjna}
\label{sec:pde13-operators}
Niech $d\geq1$, a $\Omega\subset\R^d$ będzie obszarem ograniczonym Lipschitzowskim. Wszystkie współczynniki w tym rozdziale są rzeczywiste;
analogiczne sformułowania zespolone wymagają form półtoraliniowych.
Dla funkcji $u$ rozważmy formalne wyrażenie różniczkowe
\begin{equation}\label{eq:pde13-div}
 Lu=-\sum_{i,j=1}^d\partial_i(a_{ij}\partial_j u)
       +\sum_{i=1}^d b_i\partial_i u+cu.
\end{equation}
Zakładamy $a_{ij},b_i,c\in L^\infty(\Omega)$; piszemy
$A(x)=(a_{ij}(x))$, $b(x)=(b_i(x))$. Ponadto istnieje $\lambda>0$
takie, że dla prawie każdego $x\in\Omega$ i każdego $\xi\in\R^d$
\begin{equation}\label{eq:pde13-elliptic}
 \xi^{\mathsf T}A(x)\xi\geq\lambda|\xi|^2.
\end{equation}
Warunek ten nazywamy \emph{jednostajną eliptycznością}.
Wybieramy wspólny zbiór pełnej miary, na którym nierówność
zachodzi dla wszystkich $\xi$; wystarczy jej spełnienie dla
wektorów wymiernych poza przeliczalną sumą zbiorów miary zero
i ciągłość obu stron względem $\xi$.
Niech $A_s=(A+A^{\mathsf T})/2$ oraz
$A_a=(A-A^{\mathsf T})/2$. Ponieważ $A_a^{\mathsf T}=-A_a$,
skalar $\xi^{\mathsf T}A_a\xi$ jest równy swojej transpozycji
$-\xi^{\mathsf T}A_a\xi$, a więc wynosi zero.
Zatem warunek zależy tylko od $A_s$, która jest jednostajnie
dodatnio określona; nie wymaga on symetrii $A$.
Oznaczamy przez $\|A\|_\infty$ kres istotny normy operatorowej
macierzy, a przez $\|b\|_\infty$ kres istotny normy euklidesowej
wektora. Obie liczby są skończone dzięki założeniom.

\begin{definicja}[Symbol główny i postać niedywergencyjna]
\label{def:pde13-symbol}
Symbolem głównym wyrażenia $L$ z~\eqref{eq:pde13-div} jest
$p_2(x,\xi)=\sum_{i,j}a_{ij}(x)\xi_i\xi_j$.
Formalny operator w \emph{postaci niedywergencyjnej} ma postać
$\widetilde Lu=-\sum_{i,j}a_{ij}\partial_{ij}u+
\sum_i\widetilde b_i\partial_i u+\widetilde c u$.
Jeśli $a_{ij}\in W^{1,\infty}(\Omega)$ i $u\in H^2_{\mathrm{loc}}(\Omega)$,
to lokalna reguła Leibniza dla słabych pochodnych
(twierdzenie~\ref{thm:pde9-products}, zastosowane z $p=q=2$
na kulach o domknięciu zawartym w $\Omega$, na których
$a_{ij}\in H^1$ i $\partial_j u\in H^1$) daje dla każdego $i,j$
\[
 -\partial_i(a_{ij}\partial_j u)
 =-a_{ij}\partial_{ij}u-(\partial_i a_{ij})\partial_j u
\]
w sensie dystrybucji; oba składniki po prawej są lokalnie w $L^2$.
Po zsumowaniu otrzymujemy $\widetilde L=L$ na tej klasie funkcji,
gdy $\widetilde b_j=b_j-\sum_i\partial_i a_{ij}$ oraz
$\widetilde c=c$. Przy samych współczynnikach mierzalnych
nie wolno używać tego przekształcenia jako tożsamości operatorów
na $H^1$: drugie pochodne $u$ oraz pochodne współczynników
mogą być dystrybucjami, których wskazane iloczyny nie są określone.
\end{definicja}

\begin{przyklad}[Ta sama eliptyczność, inne pojęcie rozwiązania]
\label{ex:pde13-interface}
Na $\Omega=(-1,1)$ niech $a(x)=1$ dla $x<0$ i $a(x)=4$ dla $x>0$.
Funkcja $u(x)=x$ dla $x\le0$ oraz $u(x)=x/4$ dla $x\ge0$
jest ciągła i kawałkami liniowa. Całkując przez części na obu
podprzedziałach dla dowolnej $\varphi\in C_c^\infty(-1,1)$,
otrzymujemy
\[
 -\int_{-1}^1u(x)\varphi'(x)\,\mathrm d\lambda_1(x)
 =\int_{-1}^0\varphi(x)\,\mathrm d\lambda_1(x)
 +\frac14\int_0^1\varphi(x)\,\mathrm d\lambda_1(x),
\]
gdyż człony brzegowe w zerze znikają. Zatem
$u'=\chi_{(-1,0)}+\tfrac14\chi_{(0,1)}\in L^2(-1,1)$
i $u\in H^1(-1,1)$. Ponieważ $au'=1$ prawie wszędzie,
mamy $-(au')'=0$ dystrybucyjnie.
Ponowne różniczkowanie, jak w przykładzie~\ref{ex:pde7-heaviside},
daje
\[
 \langle u'',\varphi\rangle
 =-\int_{-1}^0\varphi'(x)\,\mathrm d\lambda_1(x)
 -\frac14\int_0^1\varphi'(x)\,\mathrm d\lambda_1(x)
 =-\frac34\varphi(0).
\]
Stąd $u''=-\tfrac34\delta_0$, więc $u\notin H^2(-1,1)$.
Iloczyn $a u''$ nie jest kanonicznie określony przez klasę
$a\in L^\infty$: wybór wartości reprezentanta $a(0)$
zmieniałby formalny iloczyn z deltą, choć nie zmienia
współczynnika jako klasy funkcji. Postać niedywergencyjna
nie jest więc jedynie innym zapisem tego samego zagadnienia.
\end{przyklad}

\section{Rozwiązanie słabe i funkcjonał prawej strony}
\label{sec:pde13-weak}
Przez $H^{-1}(\Omega)$ rozumiemy przestrzeń dualną
$(H_0^1(\Omega))^*$ (definicja~\ref{def:pde11-dual}).
Niech $V=H_0^1(\Omega)$, z pełną normą Sobolewa
$\|u\|_V^2=\|u\|_{L^2}^2+\|\nabla u\|_{L^2}^2$.
W rozdziale 11 używano normy gradientu; obie normy są
równoważne na mocy wniosku~\ref{cor:pde10-Poincare},
ale tutaj normę $\|F\|_{V^*}$ liczymy względem pełnej normy.
Przestrzeń $V$ jest przestrzenią Hilberta jako domknięta
podprzestrzeń $H^1(\Omega)$
(twierdzenie~\ref{thm:pde9-banach}).
Definiujemy formę dwuliniową na $H^1(\Omega)\times H^1(\Omega)$,
w szczególności na $V\times V$,
\begin{equation}\label{eq:pde13-form}
 a(u,v)=\int_\Omega\big(A\nabla u\cdot\nabla v
 +(b\cdot\nabla u)v+cuv\big)\,\mathrm d\lambda_d(x).
\end{equation}
Z nierówności Cauchy'ego--Schwarza i ograniczoności
współczynników dla $u,v\in H^1(\Omega)$ otrzymujemy
\begin{align*}
 |a(u,v)|&\le \|A\|_\infty\|\nabla u\|_2\|\nabla v\|_2\\
 &\quad+\|b\|_\infty\|\nabla u\|_2\|v\|_2
 +\|c\|_\infty\|u\|_2\|v\|_2\\
 &\le M\|u\|_{H^1}\|v\|_{H^1},
\end{align*}
gdzie $M=\|A\|_\infty+\|b\|_\infty+\|c\|_\infty$.
Na $V$ jest to oszacowanie względem normy $\|\cdot\|_V$.
Nie zakładamy symetrii formy.

\begin{definicja}[Zagadnienie Dirichleta w sensie słabym]
\label{def:pde13-weak}
Dla $F\in V^*=H^{-1}(\Omega)$ funkcja $u\in V$ jest
\emph{słabym rozwiązaniem} zagadnienia $Lu=F$ z jednorodnym
warunkiem Dirichleta, jeżeli
\begin{equation}\label{eq:pde13-weak}
 a(u,v)=\langle F,v\rangle_{V^*,V}\qquad(v\in V).
\end{equation}
Jeśli $f\in L^2(\Omega)$, to utożsamiamy $f$ z elementem $F\in V^*$
przez $\langle F,v\rangle=\int_\Omega f(x)v(x)\,\mathrm d\lambda_d(x)$.
Z nierówności Cauchy'ego--Schwarza mamy
$|\langle F,v\rangle|\le\|f\|_2\|v\|_V$, zatem
$\|F\|_{V^*}\le\|f\|_2$. To utożsamienie jest różnowartościowe,
gdyż funkcje testowe są gęste w $L^2(\Omega)$.
Warunek brzegowy jest zakodowany w przynależności $u\in H_0^1$,
równoważnej zerowemu śladowi na mocy
twierdzenia~\ref{thm:pde10-trace-zero}.
\end{definicja}

\begin{stwierdzenie}[Zgodność ze znaczeniem dystrybucyjnym]
\label{prop:pde13-distributional}
Jeżeli $u\in H_0^1(\Omega)$ spełnia~\eqref{eq:pde13-weak}, to
$Lu=F$ jako dystrybucje na $\Omega$. Odwrotnie, gdy
$u\in H_0^1(\Omega)$ spełnia $Lu=F$ dystrybucyjnie dla
$F\in H^{-1}(\Omega)$, to spełnia~\eqref{eq:pde13-weak}.
\end{stwierdzenie}
\begin{proof}
Element $F\in V^*$ definiuje dystrybucję przez ograniczenie
do $C_c^\infty(\Omega)$: dla testów o nośniku w ustalonym
zbiorze zwartym norma $H^1$ jest ograniczona przez stałą
razy supremum modułów funkcji i jej pierwszych pochodnych.
Ograniczenie to wyznacza $F$ jednoznacznie przez gęstość testów w $V$.
Dla $\varphi\in C_c^\infty(\Omega)$ definicja pochodnej
dystrybucyjnej daje
\[
 \langle-\partial_i(a_{ij}\partial_j u),\varphi\rangle
 =\int_\Omega a_{ij}(x)(\partial_j u)(x)(\partial_i\varphi)(x)
 \,\mathrm d\lambda_d(x).
\]
Iloczyn $a_{ij}\partial_j u$ należy do $L^2$, podobnie
$b_i\partial_i u$ i $cu$, więc wszystkie wyrażenia definiują
odpowiednie dystrybucje. Po zsumowaniu po indeksach otrzymujemy
$\langle Lu,\varphi\rangle=a(u,\varphi)$.
Jeśli zachodzi~\eqref{eq:pde13-weak}, równość z $F$ zachodzi
w szczególności dla testów, co oznacza $Lu=F$ dystrybucyjnie.
Odwrotnie, wybieramy dla $v\in V$ ciąg testów $\varphi_n\to v$
w normie $V$. Z ciągłości formy i $F$ mamy
\[
 |a(u,v)-\langle F,v\rangle|
 \le (M\|u\|_V+\|F\|_{V^*})\|v-\varphi_n\|_V\longrightarrow0,
\]
gdyż równość zachodzi dla $\varphi_n$. To daje~\eqref{eq:pde13-weak}.
\end{proof}

\section{Nierówność Gårdinga i wymuszenie koercywności}
\label{sec:pde13-garding}
\begin{twierdzenie}[Elementarna nierówność Gårdinga]
\label{thm:pde13-garding}
Przy założeniach~\eqref{eq:pde13-elliptic} oraz
$a_{ij},b_i,c\in L^\infty(\Omega)$ forma~\eqref{eq:pde13-form}
spełnia
\begin{equation}\label{eq:pde13-garding}
 a(u,u)\geq\frac\lambda2\|\nabla u\|_2^2
 -\Gamma\|u\|_2^2,\qquad u\in V,
\end{equation}
gdzie można przyjąć
$\Gamma=\|b\|_\infty^2/(2\lambda)+\|c^-\|_\infty$,
$c^-=\max\{-c,0\}$. W szczególności dla
$\mu\geq\Gamma+\lambda/2$ forma
$a_\mu(u,v)=a(u,v)+\mu(u,v)_{L^2}$ jest koercywna na $V$:
$a_\mu(u,u)\geq(\lambda/2)\|u\|_V^2$.
\end{twierdzenie}
\begin{proof}
Z eliptyczności, stosowanej do $\xi=\nabla u(x)$ prawie wszędzie,
otrzymujemy
$\int_\Omega A(x)\nabla u(x)\cdot\nabla u(x)\,\mathrm d\lambda_d(x)
\geq\lambda\|\nabla u\|_2^2$.
Ponadto $c(x)u(x)^2\geq-\|c^-\|_\infty |u(x)|^2$ prawie wszędzie,
więc po scałkowaniu człon zerowego rzędu jest nie mniejszy
niż $-\|c^-\|_\infty\|u\|_2^2$.
Człon pierwszego rzędu szacujemy przez nierówność Cauchy'ego--Schwarza
i nierówność Younga $st\leq\eps s^2/2+t^2/(2\eps)$
dla $s,t\geq0$, $\eps>0$ (wynikającą z
$(\sqrt{\eps}s-t/\sqrt{\eps})^2\geq0$), wybierając $\eps=\lambda$:
\[
 \left|\int_\Omega(b(x)\cdot\nabla u(x))u(x)\,\mathrm d\lambda_d(x)\right|
 \leq \|b\|_\infty\|\nabla u\|_2\|u\|_2
 \leq\frac\lambda2\|\nabla u\|_2^2+
 \frac{\|b\|_\infty^2}{2\lambda}\|u\|_2^2.
\]
Dodanie trzech oszacowań daje~\eqref{eq:pde13-garding}.
Po dodaniu $\mu\|u\|_2^2$ otrzymujemy
$a_\mu(u,u)\geq(\lambda/2)\|\nabla u\|_2^2+
(\mu-\Gamma)\|u\|_2^2$, co jest co najmniej
$(\lambda/2)\|u\|_V^2$ przy podanym wyborze $\mu$.
\end{proof}

\begin{twierdzenie}[Istnienie dla przesuniętego operatora]
\label{thm:pde13-shift}
Przy założeniach twierdzenia~\ref{thm:pde13-garding} i dla
$\mu\geq\Gamma+\lambda/2$ dla każdego $F\in V^*$ istnieje
dokładnie jedno $u\in V$ spełniające
$a(u,v)+\mu(u,v)_{L^2}=\langle F,v\rangle$ dla wszystkich $v\in V$.
Ponadto $\|u\|_V\leq(2/\lambda)\|F\|_{V^*}$.
\end{twierdzenie}
\begin{proof}
Forma $a_\mu$ jest ograniczona ze stałą $M+\mu$, gdyż
$|(u,v)_{L^2}|\le\|u\|_V\|v\|_V$ oraz $\mu>0$.
Z poprzedniego twierdzenia jest koercywna ze stałą $\lambda/2$.
Przestrzeń $V$ jest rzeczywistą przestrzenią Hilberta, więc
twierdzenie Laxa--Milgrama~\ref{thm:pde11-LM}, które nie wymaga
symetrii formy, daje istnienie i jednoznaczność. Podstawiając $v=u$ do równania słabego,
uzyskujemy $(\lambda/2)\|u\|_V^2\leq
\langle F,u\rangle\leq\|F\|_{V^*}\|u\|_V$.
Jeżeli $u=0$, oszacowanie jest trywialne; w przeciwnym razie
dzielimy przez $\|u\|_V$.
\end{proof}

\begin{uwaga}[Dlaczego samo przesunięcie nie rozwiązuje zagadnienia?]
Wynik dotyczy operatora przesuniętego $L+\mu J:V\to V^*$,
gdzie $\langle Ju,v\rangle=(u,v)_{L^2}$, a niekoniecznie $L$.
W realizacji w $L^2$ odpowiada to przesunięciu $L+\mu I$.
Przykładowo na $(0,1)$ dla $L=-\mathrm d^2/\mathrm dx^2-\pi^2$
funkcja $\phi(x)=\sin(\pi x)$ należy do $H_0^1(0,1)$
i spełnia $L\phi=0$. Dla każdego $u\in H_0^1(0,1)$
całkowanie przez części (najpierw dla testów $u$, potem przez
gęstość) daje $a(u,\phi)=0$.
Gdyby istniało słabe rozwiązanie z prawą stroną $f=\phi$,
test $v=\phi$ dawałby $0=\|\phi\|_2^2>0$, sprzeczność.
Samo przesunięcie nie zapewnia więc rozwiązywalności
pierwotnego równania.
\end{uwaga}


\begin{uwaga}[Trzy różne wnioski z oszacowań formy]
\label{rem:pde-v42-garding}
Koercywność z twierdzenia~\ref{thm:pde11-lower-order}
pozwala bezpośrednio zastosować twierdzenie Laxa--Milgrama.
Sama nierówność Gårdinga~\eqref{eq:pde13-garding}
nie daje tego wniosku dla operatora $L$, lecz zapewnia
koercywność odpowiednio przesuniętego operatora $L+\mu I$.
Dopiero argument Fredholma w sekcji~\ref{sec:pde13-fredholm}
pozwoli wrócić do $L$, z możliwym warunkiem ortogonalności
do jądra operatora sprzężonego. Tych trzech kroków
nie należy utożsamiać.
\end{uwaga}

\section{Dirichlet: jednoznaczność, dane brzegowe i stabilność}
\label{sec:pde13-dirichlet}
\begin{twierdzenie}[Koercywność bez przesunięcia]
\label{thm:pde13-coercive}
Załóżmy, że $b=0$ i $c\geq0$ prawie wszędzie, a macierz $A$
spełnia~\eqref{eq:pde13-elliptic}. Wówczas dla każdego $F\in V^*$
zagadnienie~\eqref{eq:pde13-weak} ma dokładnie jedno rozwiązanie,
a $\|u\|_V\leq(1+C_P^2)\|F\|_{V^*}/\lambda$, gdzie
$C_P$ jest stałą nierówności Poincarégo
$\|v\|_2\leq C_P\|\nabla v\|_2$ na $V$.
Przy ogólnych $b,c$, jeśli $a(v,v)\geq\alpha\|v\|_V^2$
dla wszystkich $v\in V$ i pewnego $\alpha>0$, również istnieje
dokładnie jedno rozwiązanie, lecz oszacowanie ma wtedy postać
$\|u\|_V\leq\alpha^{-1}\|F\|_{V^*}$.
W szczególności, gdy $\Gamma C_P^2<\lambda/2$,
można przyjąć
\[
 \alpha=\frac{\lambda/2-\Gamma C_P^2}{1+C_P^2}.
\]
\end{twierdzenie}
\begin{proof}
Z wniosku~\ref{cor:pde10-Poincare} mamy
$\|v\|_V^2\leq(1+C_P^2)\|\nabla v\|_2^2$.
Dla $b=0$, $c\geq0$ eliptyczność daje
$a(v,v)\geq\lambda\|\nabla v\|_2^2
\geq\lambda(1+C_P^2)^{-1}\|v\|_V^2$.
Forma jest ograniczona na rzeczywistej przestrzeni Hilberta $V$,
więc twierdzenie Laxa--Milgrama~\ref{thm:pde11-LM} daje
istnienie i jednoznaczność.
Test $v=u$ prowadzi do
\[
 \frac{\lambda}{1+C_P^2}\|u\|_V^2
 \leq a(u,u)=\langle F,u\rangle
 \leq\|F\|_{V^*}\|u\|_V.
\]
Jeśli $u\ne0$, dzielimy przez $\|u\|_V$; dla $u=0$
oszacowanie zachodzi bez dzielenia.
Przy ogólnej koercywności stosujemy to samo twierdzenie
ze stałą $\alpha$, co daje oszacowanie $\alpha^{-1}\|F\|_{V^*}$.
W szczególnym przypadku z~\eqref{eq:pde13-garding}
i Poincarégo mamy
$a(v,v)\geq(\lambda/2-\Gamma C_P^2)\|\nabla v\|_2^2$.
Gdy współczynnik w nawiasie jest dodatni, ponowne użycie
Poincarégo daje wskazaną w tezie stałą $\alpha$.
\end{proof}

\begin{twierdzenie}[Niejednorodny warunek Dirichleta]
\label{thm:pde13-nonhom}
Niech $A,b,c$ spełniają założenia sekcji~\ref{sec:pde13-operators},
a forma $a$ będzie koercywna: $a(v,v)\geq\alpha\|v\|_V^2$
dla $v\in V$, gdzie $\alpha>0$.
Dla każdego $F\in V^*$ i $g\in H^{1/2}(\partial\Omega)$
istnieje dokładnie jedna funkcja $u\in H^1(\Omega)$ o śladzie
$\operatorname{Tr}u=g$ taka, że
$a(u,v)=\langle F,v\rangle$ dla wszystkich $v\in V$.
Ponadto
\[
 \|u\|_{H^1(\Omega)}\leq C\big(\|F\|_{V^*}
 +\|g\|_{H^{1/2}(\partial\Omega)}\big),
\]
gdzie $C$ zależy od $\alpha$, ciągłości formy i operatora
przedłużania śladu, ale nie od $F,g$.
\end{twierdzenie}
\begin{proof}
Definicja~\ref{def:pde10-halftrace} i lemat~\ref{lem:pde10-halftrace-rightinverse}
zapewniają ciągły liniowy operator
$E_{\operatorname{Tr}}:H^{1/2}(\partial\Omega)\to H^1(\Omega)$,
prawostronnie odwrotny do operatora śladu.
Kładziemy $G=E_{\operatorname{Tr}}g$; wtedy $\operatorname{Tr}G=g$
oraz $\|G\|_{H^1}\leq C_{\operatorname{Tr}}\|g\|_{H^{1/2}}$
(dla normy ilorazowej przyjętej w rozdziale 10 można wziąć
$C_{\operatorname{Tr}}=1$).
Z twierdzenia~\ref{thm:pde10-trace-zero} zbiór funkcji
o śladzie $g$ jest dokładnie $G+V$: jeśli $\operatorname{Tr}u=g$,
to $\operatorname{Tr}(u-G)=0$, a odwrotnie dodanie elementu $V$
nie zmienia śladu.
Szukamy zatem $u=G+w$, gdzie $w\in V$.
Równanie dla $w$ przyjmuje postać
$a(w,v)=\langle F,v\rangle-a(G,v)$.
Prawa strona jest ciągłym funkcjonałem na $V$ o normie
co najwyżej $\|F\|_{V^*}+M\|G\|_{H^1}$.
Z koercywności i twierdzenia Laxa--Milgrama istnieje
jedno $w$ i
$\|w\|_V\leq\alpha^{-1}
(\|F\|_{V^*}+M\|G\|_{H^1})$.
Funkcja $u=G+w$ ma ślad $g$ i przez dwuliniowość formy
spełnia równanie z tezy. Z nierówności trójkąta
\[
 \|u\|_{H^1}
 \leq \alpha^{-1}\|F\|_{V^*}
 +(1+M/\alpha)C_{\operatorname{Tr}}\|g\|_{H^{1/2}}.
\]
Można więc przyjąć
$C=\max\{\alpha^{-1},(1+M/\alpha)C_{\operatorname{Tr}}\}$.
Jeśli $u_1,u_2$ są rozwiązaniami o tych samych danych, to
$z=u_1-u_2\in V$ z twierdzenia o jądrze śladu.
Odejmując równania i testując $v=z$, dostajemy
$0=a(z,z)\geq\alpha\|z\|_V^2$, więc $z=0$.
W szczególności rozwiązanie nie zależy od wyboru funkcji $G$.
Dla dwóch różnych par danych $(F_1,g_1)$ i $(F_2,g_2)$
różnica rozwiązań spełnia to samo zagadnienie z danymi
$(F_1-F_2,g_1-g_2)$; powyższe oszacowanie daje zatem
ciągłą zależność od obu danych.
\end{proof}

\begin{lemat}[Ślad dodatniej i ujemnej części funkcji]
\label{lem:pde13-trace-parts}
Dla rzeczywistej funkcji $w\in H^1(\Omega)$ zachodzi
\[
 \operatorname{Tr}(w^+)=(\operatorname{Tr}w)^+,\qquad
 \operatorname{Tr}(w^-)=(\operatorname{Tr}w)^-
\]
w $L^2(\partial\Omega)$.
\end{lemat}
\begin{proof}
Z wniosku~\ref{cor:pde9-truncation} obie części należą do $H^1$.
Wybierzmy $w_n$ gładkie w otoczeniu $\overline\Omega$,
$w_n\to w$ w $H^1$ (wniosek~\ref{cor:pde10-density-boundary}).
Niech $\varepsilon_n=1/n$ i
$P_\varepsilon(t)=(t+\sqrt{t^2+\varepsilon^2}-\varepsilon)/2$.
Wtedy $P_\varepsilon(0)=0$, $0\le P_\varepsilon'\le1$
i $|P_\varepsilon(t)-t^+|\le\varepsilon/2$.
Funkcje $z_n=P_{\varepsilon_n}(w_n)$ są gładkie w otoczeniu
domknięcia obszaru i ograniczone w $H^1$, gdyż
$|P_\varepsilon(t)|\le|t|$ oraz
$|\nabla z_n|\le|\nabla w_n|$.
Z nierówności Lipschitzowskiej dla $t^+$ mamy
\[
 \|z_n-w^+\|_2
 \le \tfrac{\varepsilon_n}{2}\lambda_d(\Omega)^{1/2}
 +\|w_n-w\|_2\longrightarrow0.
\]
Ze słabej zwartości kul (lemat~\ref{lem:pde11-weakcompact})
wybieramy podciąg słabo zbieżny w $H^1$; silna zbieżność
w $L^2$ identyfikuje jego granicę jako $w^+$.
Ciągłość operatora śladu (twierdzenie~\ref{thm:pde10-trace})
daje słabą zbieżność śladów tego podciągu do
$\operatorname{Tr}(w^+)$.
Jednocześnie z gładkości mamy
$\operatorname{Tr}z_n=P_{\varepsilon_n}(\operatorname{Tr}w_n)$.
Ponieważ miara powierzchniowa brzegu ograniczonego obszaru
Lipschitzowskiego jest skończona, oszacowanie błędu
$P_{\varepsilon_n}$ i silna zbieżność
$\operatorname{Tr}w_n\to\operatorname{Tr}w$ w $L^2$
dają silną zbieżność $\operatorname{Tr}z_n$ do
$(\operatorname{Tr}w)^+$. Jednoznaczność granicy słabej
dowodzi pierwszej tożsamości. Drugą uzyskujemy dla $-w$.
\end{proof}

\begin{twierdzenie}[Słaba zasada minimum bez członu transportowego]
\label{thm:pde13-positivity}
Niech $b=0$, $c\geq0$ i $A$ spełnia~\eqref{eq:pde13-elliptic};
nie wymagamy symetrii macierzy $A$. Niech $u\in H^1(\Omega)$
spełnia $a(u,v)=\langle F,v\rangle$ dla wszystkich $v\in V$,
przy czym $\operatorname{Tr}u\geq0$ na $\partial\Omega$
(w sensie prawie wszędzie) i
$\langle F,v\rangle\geq0$ dla wszystkich nieujemnych $v\in V$.
Wtedy $u\geq0$ prawie wszędzie w $\Omega$.
\end{twierdzenie}
\begin{proof}
Ujemna część $u^-=\max\{-u,0\}$ należy do $H^1(\Omega)$.
Z wniosku~\ref{cor:pde9-truncation} mamy
$\nabla u^-=-\chi_{\{u<0\}}\nabla u$ prawie wszędzie.
Ponieważ ślad $u$ jest nieujemny, lemat~\ref{lem:pde13-trace-parts}
daje $\operatorname{Tr}(u^-)=0$; z twierdzenia
~\ref{thm:pde10-trace-zero} wynika $u^-\in V$.
Testujemy równanie $v=u^-$. Na zbiorze $\{u<0\}$ mamy
$u=-u^-$, $\nabla u=-\nabla u^-$, a poza nim
$u^-=\nabla u^-=0$ prawie wszędzie. Zatem
\[
 0\leq\langle F,u^-\rangle=a(u,u^-)
 =-\int_\Omega \big(A\nabla u^-\cdot\nabla u^-+c(u^-)^2\big)
 \,\mathrm d\lambda_d(x)
 \leq-\lambda\|\nabla u^-\|_2^2.
\]
Ponieważ $\lambda>0$, otrzymujemy $\|\nabla u^-\|_2=0$,
a nierówność Poincarégo na $V$ (wniosek~\ref{cor:pde10-Poincare})
daje $\|u^-\|_2=0$. Z definicji ujemnej części jest to
równoważne $u\geq0$ prawie wszędzie.
\end{proof}

\begin{twierdzenie}[Zasada porównania z dryfem o kontrolowanej dywergencji]
\label{thm:pde13-drift-comparison}
Niech $\Omega\subset\R^d$ będzie ograniczonym obszarem
Lipschitzowskim, $A\in L^\infty(\Omega;\R^{d\times d})$
spełnia~\eqref{eq:pde13-elliptic}, $b\in W^{1,\infty}(\Omega;\R^d)$,
$c\in L^\infty(\Omega)$ oraz
$c-\frac12\diver b\ge0$ prawie wszędzie. Jeżeli
$u_1,u_2\in H^1(\Omega)$ mają ślady spełniające
$\operatorname{Tr}u_1\le\operatorname{Tr}u_2$ prawie wszędzie
na $\partial\Omega$ oraz
\[
 a(u_1,\varphi)\le a(u_2,\varphi)
 \quad\text{dla wszystkich }\varphi\in H_0^1(\Omega),\ \varphi\ge0,
\]
to $u_1\le u_2$ prawie wszędzie w $\Omega$.
W szczególności założenie $\diver b=0$, $c\ge0$ zapewnia
porównanie niezależnie od wielkości normy $b$.
\end{twierdzenie}
\begin{proof}
Kładziemy $w=u_1-u_2$ i $v=w^+$.
Wniosek~\ref{cor:pde9-truncation} daje $v\in H^1$
oraz $\nabla v=\chi_{\{w>0\}}\nabla w$.
Z liniowości śladu i założenia $\operatorname{Tr}w\le0$,
więc lemat~\ref{lem:pde13-trace-parts} daje $\operatorname{Tr}v=0$.
Stąd $v\in H_0^1(\Omega)$ na mocy twierdzenia o jądrze śladu.
Ponieważ $v\ge0$, wolno wybrać $\varphi=v$, co daje
$a(w,v)\le0$.
Człon pierwszego rzędu obliczamy bez założenia małości $b$:
\[
 \int_\Omega (b\cdot\nabla w)v\,\mathrm d\lambda_d(x)
 =\frac12\int_\Omega b\cdot\nabla(v^2)\,\mathrm d\lambda_d(x)
 =-\frac12\int_\Omega(\diver b)v^2\,\mathrm d\lambda_d(x).
\]
Pierwsza równość wynika z $\nabla w=\nabla v$ na $\{w>0\}$
i $v=0$ poza tym zbiorem oraz z reguły Leibniza
(twierdzenie~\ref{thm:pde9-products}): $v^2\in W^{1,1}$
i $\nabla(v^2)=2v\nabla v$.
Dla ostatniej równości wybieramy $v_n\in C_c^\infty(\Omega)$
zbieżne do $v$ w $H^1$. Mamy
\begin{align*}
 \|v_n^2-v^2\|_1
 &\le\|v_n-v\|_2\|v_n+v\|_2,\\
 \|\nabla(v_n^2)-\nabla(v^2)\|_1
 &\le 2\|v_n-v\|_2\|\nabla v_n\|_2
 +2\|v\|_2\|\nabla v_n-\nabla v\|_2.
\end{align*}
Obie prawe strony dążą do zera, zatem $v_n^2\to v^2$
w $W^{1,1}$. Dla testu $v_n^2$ definicja słabej dywergencji
daje
$\int_\Omega b\cdot\nabla(v_n^2)\,\mathrm d\lambda_d(x)
=-\int_\Omega(\diver b)v_n^2\,\mathrm d\lambda_d(x)$.
Ponieważ $b,\diver b$ są ograniczone, możemy przejść
do granicy w obu całkach. W ten sposób uzyskujemy równość
bez członu brzegowego, nie zakładając $v^2\in H^1$.
Wobec tego
\[
 0\ge a(w,v)=\int_\Omega\left(A\nabla v\cdot\nabla v+
 (c-\tfrac12\diver b)v^2\right)\,\mathrm d\lambda_d(x)
 \ge\lambda\|\nabla v\|_2^2.
\]
Nierówność Poincarégo daje $\|v\|_2=0$, czyli
$(u_1-u_2)^+=0$ prawie wszędzie. Otrzymujemy żądane porównanie.
Jest to argument energetyczny, a nie twierdzenie o punktowej
regularności rozwiązań. W szczególności dla dwóch rozwiązań
o tych samych danych zastosowanie porównania w obu kierunkach
daje jednoznaczność.

Przy tych samych założeniach znakowych forma jest również
koercywna na $V$. Istotnie, dla dowolnego $z\in V$ argument
aproksymacji kwadratów zastosowany do $z$ daje
\[
 a(z,z)=\int_\Omega
 \big(A\nabla z\cdot\nabla z+(c-\tfrac12\diver b)z^2\big)
 \,\mathrm d\lambda_d(x)
 \ge\lambda\|\nabla z\|_2^2.
\]
Z Poincarégo otrzymujemy koercywność ze stałą
$\lambda/(1+C_P^2)$, bez warunku małości $b$.
Twierdzenia~\ref{thm:pde13-coercive} i~\ref{thm:pde13-nonhom}
zapewniają więc także istnienie przy jednorodnych i niejednorodnych
danych Dirichleta.
\end{proof}



\section{Alternatywa Fredholma dla eliptycznych zagadnień Dirichleta}
\label{sec:pde13-fredholm}
Dla współczynników dolnego rzędu bez warunku znakowego forma $a$
może nie być koercywna. Mimo to operator $L:V\to V^*$,
$\langle Lu,v\rangle=a(u,v)$, ma precyzyjną własność Fredholma.
Nie należy go mylić z realizacją $L$ w $L^2$, której dziedzina
wymaga dodatkowej definicji.

\begin{twierdzenie}[Alternatywa Fredholma dla $L:V\to V^*$]
\label{thm:pde13-fredholm}
Niech $\Omega\subset\R^d$ będzie obszarem ograniczonym Lipschitzowskim, a $A,b,c\in L^\infty$ niech spełniają
jednostajną eliptyczność~\eqref{eq:pde13-elliptic}.
Wówczas $L:V\to V^*$ jest operatorem Fredholma indeksu $0$:
ma domknięty obraz, skończeniewymiarowe jądro i skończony
kowymiar obrazu, równy wymiarowi jądra.
Definiując operator transponowany $L^{\mathrm t}:V\to V^*$ przez
$\langle L^{\mathrm t}z,v\rangle=a(v,z)$, mamy
$\dim\ker L=\dim\ker L^{\mathrm t}<\infty$.
Równanie $Lu=F$, $F\in V^*$, ma rozwiązanie wtedy i tylko wtedy, gdy
\begin{equation}\label{eq:pde13-compat}
 \langle F,z\rangle=0\qquad(z\in\ker L^{\mathrm t}).
\end{equation}
Rozwiązania różnią się o elementy $\ker L$. W szczególności
gdy $\ker L=\{0\}$, dla każdego $F\in V^*$ istnieje jedno rozwiązanie.
\end{twierdzenie}
\begin{proof}
\textbf{Krok 1: odwracalny operator przesunięty.}
Wybieramy $\mu\geq\Gamma+\lambda/2$.
Z twierdzenia~\ref{thm:pde13-shift} operator
$L_\mu=L+\mu J:V\to V^*$ jest izomorfizmem ograniczonym,
przy czym $J:V\to V^*$ definiujemy wzorem
$\langle Ju,v\rangle=\int_\Omega uv\,\mathrm d\lambda_d(x)$.

\textbf{Krok 2: zwartość zaburzenia.}
Zanurzenie $i:V\hookrightarrow L^2(\Omega)$ jest zwarte
na mocy twierdzenia Rellicha--Kondrachova. Operator
$j:L^2(\Omega)\to V^*$, $\langle jf,v\rangle=(f,v)_{L^2}$,
jest ograniczony, ponieważ $\|v\|_2\leq\|v\|_V$.
Zatem $J=j\circ i$ jest zwarty jako złożenie operatora
zwartego z ograniczonym. W konsekwencji operator
$K=\mu L_\mu^{-1}J:V\to V$ jest zwarty.

\textbf{Krok 3: sprowadzenie do alternatywy dla operatora zwartego.}
Mamy równość operatorów
$L=L_\mu(I-K)$; istotnie
$L_\mu K=\mu J$.
Twierdzenie~\ref{thm:pde12-general-fredholm} zostało
udowodnione dla zespolonych przestrzeni Hilberta;
jego dowód oparty na aproksymacji operatorami skończonego rzędu
pozostaje bez zmian prawdziwy w rzeczywistych przestrzeniach
Hilberta. Daje to, że $I-K:V\to V$ ma domknięty obraz,
skończeniewymiarowe jądro i kowymiar równy wymiarowi jądra.
Ponieważ $L_\mu$ jest izomorfizmem, własności te przechodzą
na $L:V\to V^*$, a zatem $\operatorname{ind}L=0$.

\textbf{Krok 4: dokładny warunek zgodności.}
Dla $z\in V$ warunek $z\in\ker L^{\mathrm t}$ oznacza
$a(v,z)=0$ dla wszystkich $v\in V$. To jest równoważne temu,
że ciągły funkcjonał na $V^*$ zadany przez
$F\mapsto\langle F,z\rangle$ znika na $\operatorname{Ran}L$.
Każdy ciągły funkcjonał na $V^*$ ma tę postać, ponieważ
$V$ jest refleksywna (jako przestrzeń Hilberta).
Przestrzeń funkcjonałów zerujących obraz $L$ ma zatem
wymiar $\dim\ker L^{\mathrm t}$.
Ponieważ obraz jest domknięty, twierdzenie Hahna--Banacha
lub identyfikacja przestrzeni ilorazowej daje
$\operatorname{Ran}L=\{F\in V^*: \langle F,z\rangle=0
\text{ dla każdego } z\in\ker L^{\mathrm t}\}$.
Kowymiar obrazu jest równy $\dim\ker L^{\mathrm t}$,
a z kroku 3 także $\dim\ker L$. Wreszcie różnica rozwiązań
jednego równania należy do jądra $L$.
\end{proof}

\begin{przyklad}[Zgodność przy rezonansie]
\label{ex:pde13-resonance}
Dla $\Omega=(0,1)$ i $L=-\mathrm d^2/\mathrm dx^2-\pi^2$
z warunkiem Dirichleta mamy
$\ker L=\ker L^{\mathrm t}=\operatorname{span}\{\sin(\pi x)\}$.
Równanie $-u''-\pi^2u=f$, $u(0)=u(1)=0$,
ma słabe rozwiązanie dokładnie wtedy, gdy
$\int_0^1 f(x)\sin(\pi x)\,\mathrm d\lambda_1(x)=0$.
Gdy ten warunek zachodzi, rozwiązanie jest wyznaczone
z dokładnością do dodania wielokrotności $\sin(\pi x)$.
W szczególności dla $f=\sin(\pi x)$ rozwiązanie nie istnieje.
\end{przyklad}

\section{Zagadnienie Neumanna: warunek zgodności i stałe}
\label{sec:pde13-neumann}
Dla $b=c=0$ i jednorodnej eliptyczności rozważmy
$-\diver(A\nabla u)=f$ wraz z zadanym strumieniem na brzegu.
Przy słabych współczynnikach nie musimy definiować strumienia
punktowo; dane brzegowe można wprowadzić przez funkcjonał śladu.

\begin{lemat}[Nierówność Poincarégo--Wirtingera]
\label{lem:pde13-wirtinger}
Jeżeli $\Omega$ jest obszarem ograniczonym Lipschitzowskim,
to istnieje $C_W>0$ takie, że dla każdego $u\in H^1(\Omega)$
\[
\|u-u_\Omega\|_{L^2(\Omega)}\le C_W\|\nabla u\|_{L^2(\Omega)},
\qquad u_\Omega=\frac1{\lambda_d(\Omega)}
\int_\Omega u\,\mathrm d\lambda_d(x).
\]
\end{lemat}
\begin{proof}
Gdyby twierdzenie było fałszywe, znaleźlibyśmy $v_n\in H^1$
o średniej zero takie, że $\|v_n\|_2=1$ i
$\|\nabla v_n\|_2\to0$ (odejmujemy średnią i normalizujemy).
Ciąg $v_n$ jest ograniczony w $H^1$, więc dzięki
zwartości $H^1(\Omega)\hookrightarrow L^2(\Omega)$
ma podciąg zbieżny silnie w $L^2$ do $v$. Po wybraniu
podciągu słabo zbieżnego w $H^1$ otrzymujemy
$\nabla v=0$ w sensie dystrybucyjnym.
Ponieważ $\Omega$ jest spójna, $v$ jest prawie wszędzie stała.
Ciągłość funkcjonału $w\mapsto\int_\Omega w$ oraz średnie
zerowe pokazują, że ta stała wynosi zero.
Ale silna zbieżność w $L^2$ wymusza $\|v\|_2=1$,
co jest sprzecznością.
\end{proof}

\begin{twierdzenie}[Neumann: istnienie modulo stałe]
\label{thm:pde13-neumann}
Niech $\Omega$ będzie obszarem ograniczonym Lipschitzowskim, a $A\in L^\infty(\Omega;\R^{d\times d})$
spełnia~\eqref{eq:pde13-elliptic}. Niech
$f\in L^2(\Omega)$ i $h\in L^2(\partial\Omega)$,
przy czym na brzegu całkujemy względem miary powierzchniowej
$\sigma=\mathcal H^{d-1}|_{\partial\Omega}$.
Istnieje $u\in H^1(\Omega)$ spełniające
\begin{equation}\label{eq:pde13-neumann}
\int_\Omega A\nabla u\cdot\nabla v\,\mathrm d\lambda_d(x)
=\int_\Omega fv\,\mathrm d\lambda_d(x)
+\int_{\partial\Omega} h\,\operatorname{Tr}v\,\mathrm d\sigma
\quad(v\in H^1(\Omega))
\end{equation}
wtedy i tylko wtedy, gdy
\begin{equation}\label{eq:pde13-neumann-comp}
\int_\Omega f\,\mathrm d\lambda_d(x)
+\int_{\partial\Omega}h\,\mathrm d\sigma=0.
\end{equation}
Rozwiązanie jest jednoznaczne modulo stałe. Istnieje dokładnie
jedno rozwiązanie o średniej zero; spełnia ono
$\|u\|_{H^1}\leq C(\|f\|_2+\|h\|_{L^2(\partial\Omega)})$.
\end{twierdzenie}
\begin{proof}
\textbf{Konieczność.} Podstawiamy $v=1$ do
\eqref{eq:pde13-neumann}; gradient stałej znika i
otrzymujemy dokładnie~\eqref{eq:pde13-neumann-comp}.

\textbf{Wystarczalność.} Rozważmy domkniętą podprzestrzeń
$V_\diamond=\{v\in H^1(\Omega):v_\Omega=0\}$.
Na niej norma $\|\nabla v\|_2$ jest równoważna normie
$H^1$ na mocy lematu~\ref{lem:pde13-wirtinger}.
Forma $B(u,v)=\int_\Omega A\nabla u\cdot\nabla v
\,\mathrm d\lambda_d(x)$ jest ciągła oraz koercywna
względem tej normy, ponieważ
$B(v,v)\geq\lambda\|\nabla v\|_2^2$.
Funkcjonał po prawej stronie~\eqref{eq:pde13-neumann}
jest ciągły na $V_\diamond$: przez nierówność Höldera
oraz ciągłość śladu jego moduł jest nie większy niż
$C(\|f\|_2+\|h\|_{L^2(\partial\Omega)})\|v\|_{H^1}$.
Z twierdzenia Laxa--Milgrama istnieje jedno $u\in V_\diamond$
spełniające równanie dla wszystkich $v\in V_\diamond$.
Dla dowolnego $v\in H^1$ piszemy $v=(v-v_\Omega)+v_\Omega$.
Pierwsza część jest w $V_\diamond$, a prawa strona na stałej
znika dzięki zgodności~\eqref{eq:pde13-neumann-comp}.
Równanie jest zatem prawdziwe dla całego $H^1$.
Oszacowanie wynika z Laxa--Milgrama i lematu Wirtingera.
Jeśli $u_1,u_2$ są rozwiązaniami, ich różnica $w$ spełnia
$B(w,v)=0$ dla wszystkich $v\in H^1$. Biorąc $v=w$,
otrzymujemy $\nabla w=0$, a spójność $\Omega$ implikuje,
że $w$ jest prawie wszędzie stała.
\end{proof}

\begin{uwaga}[Znak strumienia]
Dla gładkiego $u$ i $A$ całkowanie przez części daje
$\int_\Omega A\nabla u\cdot\nabla v
=\int_\Omega[-\diver(A\nabla u)]v
+\int_{\partial\Omega}(A\nabla u\cdot\nu)v$.
Dlatego $h$ w~\eqref{eq:pde13-neumann} oznacza
\emph{zewnętrzny} strumień $A\nabla u\cdot\nu$.
Przy danych szorstkich jest to jedynie interpretacja
funkcjonałowa; nie zakładamy istnienia punktowego śladu
wektora $A\nabla u$.
\end{uwaga}

\section{Robin: człon brzegowy przywracający jednoznaczność}
\label{sec:pde13-robin}
\begin{lemat}[Nierówność Poincarégo z wyrazem brzegowym]
\label{lem:pde13-boundary-poincare}
Na obszarze ograniczonym Lipschitzowskim istnieje
$C>0$ takie, że dla wszystkich $u\in H^1(\Omega)$
\[
\|u\|_{H^1}^2\le C\big(\|\nabla u\|_2^2+
 \|\operatorname{Tr}u\|_{L^2(\partial\Omega)}^2\big).
\]
\end{lemat}
\begin{proof}
Gdyby oszacowanie nie zachodziło, mielibyśmy ciąg
$\|u_n\|_{H^1}=1$ i $\|\nabla u_n\|_2^2+
\|\operatorname{Tr}u_n\|_{L^2(\partial\Omega)}^2\to0$.
Ze zwartości zanurzenia $H^1\hookrightarrow L^2$ i słabej
zwartości kuli w $H^1$ otrzymujemy podciąg słabo zbieżny
w $H^1$, silnie w $L^2$ do pewnego $u$.
Ponieważ $\nabla u_n\to0$ w $L^2$, mamy $\nabla u=0$,
a więc $u$ jest stała. Z ciągłości śladu $H^1\to
L^2(\partial\Omega)$ otrzymujemy słabą zbieżność śladów
do $\operatorname{Tr}u$, lecz ślady dążą silnie do zera.
Zatem stała $u$ ma zerowy ślad, czyli $u=0$.
Wówczas $\|u_n\|_2\to0$ oraz $\|\nabla u_n\|_2\to0$,
wbrew $\|u_n\|_{H^1}=1$.
\end{proof}

\begin{twierdzenie}[Zagadnienie Robina]
\label{thm:pde13-robin}
Niech $\Omega$ będzie obszarem ograniczonym Lipschitzowskim, $A\in L^\infty$ spełnia
\eqref{eq:pde13-elliptic}, a $\beta\in L^\infty(\partial\Omega)$
spełnia $\beta\geq\beta_0>0$ prawie wszędzie.
Dla każdego $f\in L^2(\Omega)$ i $h\in L^2(\partial\Omega)$
istnieje dokładnie jedno $u\in H^1(\Omega)$ spełniające
\begin{equation}\label{eq:pde13-robin}
\int_\Omega A\nabla u\cdot\nabla v\,\mathrm d\lambda_d(x)
+\int_{\partial\Omega}\beta\,\operatorname{Tr}u\,
\operatorname{Tr}v\,\mathrm d\sigma
=\int_\Omega fv\,\mathrm d\lambda_d(x)
+\int_{\partial\Omega}h\,\operatorname{Tr}v\,\mathrm d\sigma
\end{equation}
dla wszystkich $v\in H^1(\Omega)$, a ponadto
$\|u\|_{H^1}\leq C(\|f\|_2+\|h\|_{L^2(\partial\Omega)})$.
Formalnie warunek brzegowy ma postać
$A\nabla u\cdot\nu+\beta u=h$.
\end{twierdzenie}
\begin{proof}
Ciągłość formy w lewej stronie wynika z ograniczoności $A,\beta$
oraz ciągłości operatora śladu $H^1\to L^2(\partial\Omega)$.
Jej wartość dla $v=u$ jest nie mniejsza niż
$\lambda\|\nabla u\|_2^2+
\beta_0\|\operatorname{Tr}u\|_{L^2(\partial\Omega)}^2$,
a z lematu~\ref{lem:pde13-boundary-poincare} wynika koercywność
na całym $H^1$. Prawa strona jest ciągłym funkcjonałem
na $H^1$ przez Höldera i twierdzenie o śladzie.
Wszystkie założenia Laxa--Milgrama są więc spełnione;
otrzymujemy istnienie, jednoznaczność i oszacowanie.
\end{proof}

\section{Przykład obliczeniowy: strumień przez granicę ośrodków}
\label{sec:pde13-transmission}
Niech $\Omega=(-1,1)$, $f=1$ i
$a(x)=1$ dla $x<0$, $a(x)=4$ dla $x>0$.
Rozważamy $-(a u')'=1$ wraz z $u(-1)=u(1)=0$.
Na obu podprzedziałach $u''=-1/a$, a zatem
\begin{equation}\label{eq:pde13-exact}
 u(x)=\begin{cases}
 -x^2/2-3x/10+1/5,& -1\le x\le0,\\
 -x^2/8-3x/40+1/5,& 0\le x\le1.
 \end{cases}
\end{equation}
\begin{center}
\begin{tikzpicture}[x=3.1cm,y=0.7cm,>=Stealth]
  \fill[blue!8] (-1,0) rectangle (0,1.35);
  \fill[orange!13] (0,0) rectangle (1,1.35);
  \draw[thick] (-1,0) rectangle (1,1.35);
  \draw[thick,dashed] (0,0)--(0,1.35);
  \node at (-.52,.87) {$a_-=1$};
  \node at (.52,.87) {$a_+=4$};
  \node at (-.52,.38) {ośrodek I};
  \node at (.52,.38) {ośrodek II};
  \draw[->] (-1.15,-.24)--(1.18,-.24) node[right] {$x$};
  \foreach \x/\lab in {-1/-1,0/0,1/1}{
    \draw (\x,-.31)--(\x,-.17);\node[below] at (\x,-.31) {$\lab$};}
  \node[above] at (0,1.35) {interfejs};
\end{tikzpicture}

\small\emph{Dwa ośrodki o różnych współczynnikach: ciągłe $u$ i $a u'$, lecz skok $u'$.}
\end{center}

Wartość $u(0^-)=u(0^+)=1/5$ jest ciągła, natomiast
$u'(0^-)=-3/10$ i $u'(0^+)=-3/40$. Mimo skoku pochodnej
strumień $a(x)u'(x)=-x-3/10$ jest ciągły w $0$.

\begin{stwierdzenie}[Warunki transmisyjne z równania słabego]
\label{prop:pde13-transmission}
Niech $a$ jest dodatnią stałą $a_-$ na $(-1,0)$ i
$a_+$ na $(0,1)$, $f\in L^2(-1,1)$, a $u\in H^1(-1,1)$
ma na obu podprzedziałach regularność $H^2$.
Wówczas $-(au')'=f$ w $\mathcal D'(-1,1)$ wtedy i tylko wtedy,
gdy $-a_\pm u''=f$ prawie wszędzie po odpowiednich stronach,
a ślady strumienia spełniają
$a_-u'(0^-)=a_+u'(0^+)$.
\end{stwierdzenie}
\begin{proof}
Z ciągłości funkcji z $H^1(-1,1)$ w wymiarze jednym
wynika $u(0^-)=u(0^+)$. Testując równanie funkcjami o nośniku
w każdym z podprzedziałów, dostajemy tam równania klasyczne
prawie wszędzie. Dla dowolnej $\varphi\in C_c^\infty(-1,1)$
całkujemy osobno po obu przedziałach:
\[
\int_{-1}^1 au'\varphi'\,\mathrm d\lambda_1
=\int_{-1}^1 f\varphi\,\mathrm d\lambda_1
 +\big(a_-u'(0^-)-a_+u'(0^+)\big)\varphi(0).
\]
Zatem równość dystrybucyjna na całym przedziale jest
równoważna znikaniu współczynnika przy $\varphi(0)$.
W drugą stronę te same rachunki, przy zerowym skoku strumienia,
dają równanie globalne. Dla~\eqref{eq:pde13-exact} skok strumienia
wynosi zero, co potwierdza poprawność rozwiązania.
\end{proof}

\section{Zakres otrzymanych rezultatów}
\label{sec:pde13-summary}
Jednostajna eliptyczność zapewnia kontrolę części gradientowej,
lecz człony dolnego rzędu mogą zniszczyć koercywność.
Nierówność Gårdinga pozwala odzyskać ją po przesunięciu,
a zwartość zanurzenia Sobolewa prowadzi do alternatywy Fredholma
dla problemu pierwotnego. W zagadnieniu Neumanna przeszkodą
jest przestrzeń stałych, zaś dodatni człon brzegowy w problemie
Robina usuwa tę niejednoznaczność.

Wszystkie twierdzenia o istnieniu dotyczą rozwiązań słabych.
Nie twierdzimy, że sam warunek eliptyczności i ograniczoność
mierzalnych współczynników dają $H^2$ lub gładkość klasyczną.
Przykład transmisyjny pokazuje dokładnie, gdzie taka
regularność może zawieść. Regularność wewnętrzną i brzegową
rozwiniemy w kolejnych rozdziałach.

% Iteracja 14: wykład, do włączenia przed dodatkiem z kodem.
